%%%%%%%%%%%%%%%%%%%%%%%%%%%%%%%%%%%%%%%%%
% Medium Length Professional CV
% LaTeX Template
% Version 2.0 (8/5/13)
%
% This template has been downloaded from:
% http://www.LaTeXTemplates.com
%
% Original author:
% Rishi Shah 
%
% Important note:
% This template requires the resume.cls file to be in the same directory as the
% .tex file. The resume.cls file provides the resume style used for structuring the
% document.
%
%%%%%%%%%%%%%%%%%%%%%%%%%%%%%%%%%%%%%%%%%

%----------------------------------------------------------------------------------------
%	PACKAGES AND OTHER DOCUMENT CONFIGURATIONS
%----------------------------------------------------------------------------------------
\newcommand{\kangda}[1]{\textcolor{brown}{\small{\bf [ #1 -- Kangda ]}}}
\documentclass[a4paper, 12pt]{resume} % Use the custom resume.cls style
\usepackage{url}
\usepackage[left=0.6in,top=0.6in,right=0.6in,bottom=0.6in]{geometry} % Document margins
\newcommand{\tab}[1]{\hspace{.1\textwidth}\rlap{#1}}
\newcommand{\itab}[1]{\hspace{0em}\rlap{#1}}
\name{Kangda Wei} % Your name
\address{L.F. Peterson Building, Room 330, 435 Nagle St, College Station, TX 77843} % Your address
%\address{123 Pleasant Lane \\ City, State 12345} % Your secondary addess (optional)
\address{\url{https://weikangda.github.io/kangda.github.io}}
\address{(919) 904-5280 \\ kangda@tamu.edu \\ hansonwei666@gmail.com} % Your phone number and email

\begin{document}

%----------------------------------------------------------------------------------------
%	Research Interest SECTION
%----------------------------------------------------------------------------------------
\begin{rSection}{Research Interest}
Natural Language Processing, Large Language Model, LLM Reasoning, Post-Training, Multi-Modal, Agentic Systems
\end{rSection}

%----------------------------------------------------------------------------------------
%	Education Background SECTION
%----------------------------------------------------------------------------------------
\begin{rSection}{Education}
{\bf Texas A\&M University} \hfill {College Station, TX \\
{\em Doctor of Philosophy in Computer Science. Advisor: Ruihong Huang.} \hfill Aug 2023 - May 2028 (Expected)} \\
%\begin{itemize}[leftmargin=*]
%\item {\bf Graduate Research Assistant}: Advised by Professor Ruihong Huang.
%\end{itemize}
{\bf University of North Carolina at Chapel-Hill} \hfill {Chapel-Hill, NC \\
{\em Master of Science in Computer Science. Advisor: Shashank Srivastava.} \hfill Aug 2022 - May 2023} \\
%\begin{itemize}[leftmargin=*]
%\item {\bf Graduate Research Assistant}: Advised by Professor Shashank Srivastava.
%\item {\bf Graduate Teaching Assistant}: COMP431 Internet Service and Protocol, COMP211 System Fundamentals
%\end{itemize}
{\bf University of North Carolina at Chapel-Hill} \hfill {Chapel-Hill, NC \\
{\em B.S. in Computer Science, B.S. in Statistics and Operational Research, GPA: 3.812/4.0} \hfill Aug 2019 - May 2022}
%\begin{itemize}[leftmargin=*]
%\item {\bf GPA}: 3.812/4.0
%\item {\bf Relevant coursework}: Foundation of Programming, Data Structures, Algorithms and Analysis, Computational Photograph, Optimization, Structured Prediction, Machine Learning \& Deep Learning, Programming Language, Time Series, Simulation Analysis, Stochastic Modeling, Linear Algebra, Computer Organization, Reinforcement Learning
%\end{itemize}
%{\bf Online courses}
%\begin{itemize}[leftmargin=*]
%\item Stanford CS224n (NLP), Machine Learning, Deep Learning Specialization Series taught by Andrew Ng on Coursera

%\end{itemize}
\end{rSection}

%----------------------------------------------------------------------------------------
%	WORK EXPERIENCE SECTION
%----------------------------------------------------------------------------------------
\begin{rSection}{Industry Experience}
{\bf Applied Scientist Intern} \\
{\em Amazon}\hfill{Santa Clara, CA}\\
{\em Self-Attribution-Guided Self-Evolving, Mentor: Hao Chen, Manager: Gokce Keskin}\hfill{May 2026 - Dec 2026} \\
{\bf Research Scientist Intern} \\
{Bosch Center for Artificial Intelligence (BCAI)}\hfill{Sunnyvale, CA}\\
{\em Multi-Agent Video Understanding, Mentor: Zhengyu Zhou}\hfill{May 2024 - Aug 2024}
\end{rSection}

%----------------------------------------------------------------------------------------
%	Research EXPERIENCE SECTION
%----------------------------------------------------------------------------------------

\begin{rSection}{Research Experience}
{\bf OpenRSI-AI-Infra-Bench: Benchmarking Recursive Self-Improvement of Agents on AI Serving}
\begin{itemize}[noitemsep, topsep=0pt, leftmargin=*]
\item Co-led (co-first author) a benchmark, built with the SGLang-Omni team, where coding agents iteratively optimize a production serving stack on one GPU; each speedup is verified under quality and latency constraints and becomes the next baseline.
\item Built 8 reproducible serving optimization tasks for open-weight TTS/ASR models (e.g., Qwen3-TTS, Whisper large-v3, CosyVoice3) with pinned code, dependencies, and evaluation protocols.
\item Benchmarked Claude Code and Codex in 24-hour runs on one H100, reaching up to $6.3\times$ verified throughput; agent-found optimizations led to 14 PRs merged into upstream SGLang-Omni.
\item \emph{Status:} Paper under preparation for submission to ICML 2027; part of the OpenRSI-Index Project (co-author).
\end{itemize}

{\bf Self-Attribution-Guided Prompt Evolution for  Large Language Models}
\begin{itemize}[noitemsep, topsep=0pt, leftmargin=*]
\item Proved that prompt optimizers steering only by task reward are provably blind among tied candidates, and that a prefill-only teacher-forced log-likelihood, the self-attribution score, certifies a gradient-free lower bound on the information they discard.
\item Developed SAGE, a drop-in extension to a state-of-the-art reflective prompt optimizer over a frozen model, wiring this signal into parent selection and child acceptance as an accuracy-gated constraint rather than a reward.
\item Achieved the highest Accuracy@4 on four open-weights models spanning a $30\times$ size range and six tasks, beating the strongest accuracy-only optimizer by up to $+4.0$ points while cutting regressions $4\times$, with gains holding under difficulty shift, cross-model transfer, and black-box deployment.
\item \emph{Status:} Paper submitted to ICLR 2027; conducted during Applied Scientist internship at Amazon.
\end{itemize}


{\bf Accelerating Mathematical Reasoning Model Training via Diversity-Aware Reward Reweighting}
\begin{itemize}[noitemsep, topsep=0pt, leftmargin=*]
\item Developed MMR-GRPO, a training-efficient approach that integrates Maximal Marginal Relevance into Group Relative Policy Optimization to accelerate convergence by prioritizing semantically diverse completions and downweighting redundant solutions.
\item Designed a parameter-free adaptive mechanism that automatically balances diversity and reward quality without hyperparameter tuning, enabling robust performance across different training configurations.
\item Achieved 47.9\% reduction in training steps and 70.2\% reduction in wall-clock training time while maintaining comparable peak performance across three model sizes (1.5B, 7B, 8B), three RL methods (GRPO, DR-GRPO, DAPO), and five mathematical reasoning benchmarks (AIME 2024, AMC 2023, MATH-500, Minerva Math, OlympiadBench).
\item Demonstrated minimal computational overhead (1-5\% per-step) through vectorized GPU operations for embedding extraction and similarity computation, making the approach practical for resource-constrained settings.
%\item \emph{Status:} Paper submitted to ACL 2026; code and models to be released upon acceptance.
\end{itemize}

{\bf Mitigating Bias and Improving Reasoning in LLMs via Exploratory Thinking and RL}
%{\em Texas A\&M University}\hfill{College Station, TX}\\
%{\em Research Assistant, Mentor: Prof.Ruihong Huang}\hfill{Feb 2025 - Aug 2025}
\begin{itemize}[noitemsep, topsep=0pt, leftmargin=*]
\item Proposed a novel data generation framework that encourages exploratory thinking in LLMs by prompting them to reason over structurally identical, morally ambiguous story pairs with male and female protagonists. 
\item Designed a guided judgement reconciliation process to uncover and mitigate biased moral reasoning, producing balanced gender-neutral story-judgement pairs. 
\item Performed SFT and DPO on LLMs, achieving substantial gender bias reduction on WinoBias and GenMO benchmarks without sacrificing performance on MMLU and TruthfulQA.
%\item \emph{Ongoing:} Extending framework to other forms of social bias (e.g., race, social status) 
\end{itemize}


%{\bf Multi-Agent Systems with LLMs/VLMs for Nuclear Safety Anlysis and Video Understanding}
%{\em Texas A\&M University; Robert Bosch Research and Technology Center North America} \hfill{College Station, TX \& Sunnyvale, CA}\\
%{\em Research Assistant (Mentor: Prof. Ruihong Huang); NLP \& LLMs Intern (Mentor: Zhengyu Zhou)} \hfill{May 2024 -- Present}
%\begin{itemize}[noitemsep, topsep=0pt,leftmargin=*]
%\item Designed and implemented multi-agent systems that leverage LLMs and VLMs for two domains: 
%  \begin{itemize}[noitemsep, topsep=0pt,leftmargin=*]
%    \item {\em Nuclear Safety:} Developed a BFS-based multi-agent framework for automatic nuclear fault tree generation in collaboration with Idaho National Laboratory (INL). Built a large-scale nuclear safety dataset via OCR and expert validation; benchmarked multiple LLMs/VLMs for fault analysis tasks. 
%    \item {\em Video Understanding:} Proposed an end-to-end multi-agent multi-modal framework for advanced indexing of lecture/presentation-style videos. Combined VLM-based shot detection with multi-agent collaboration to capture rich cross-modal information for scalable downstream use. Conducted systematic evaluations, including intrinsic human evaluation (Amazon Mechanical Turk) and extrinsic QA benchmarks, on both public and enterprise datasets. 
%  \end{itemize}
%\item Outcomes: one paper published/submitted to top-tier NLP conferences (EMNLP/ACL/AACL 2025) and a patent application filed with the USPTO. 
%\end{itemize}

%{\bf Event Understanding and Coreference in Long Documents}
%{\em Texas A\&M University}\hfill{College Station, TX}\\
%{\em Research Assistant, Mentor: Prof. Ruihong Huang}\hfill{Aug 2023 -- Feb 2025}
%\begin{itemize}[noitemsep, topsep=0pt,leftmargin=*]
%\item Introduced {\em LegalCore}, the first dataset in the legal domain annotated with comprehensive event and event coreference information. 
%\item Benchmarked mainstream LLMs on discourse-level event relation extraction and event coreference tasks, showing that both remain challenging for LLMs. 
%\item Provided strong supervised baselines by fine-tuning smaller language models and conducted extensive empirical analysis, highlighting weaknesses of GPT models in hallucination control, transitivity, and long-range dependency capture.  
%\item Published in {\em Findings of EMNLP 2024} and {\em ACL 2025}.  
%\end{itemize}

%{\bf When Do Decompositions Help for Machine Reading?} \\
%{\em  Johns Hopkins University}\hfill{Baltimore, MD}\\
%{\em Visiting Research Assistant, Mentor: Prof.Benjamin van Durme}\hfill{May 2022 - May 2023}
%\begin{itemize}[noitemsep, topsep=0pt,leftmargin=*]
%\item Explored the effect of decomposition on machine reading with an exhaustive set of variants across a range of models over the high-level Question Decomposition Meaning Representation (QDMR) BREAK dataset using various LLMs, including T-5, Alpaca, and LLaMA.
%\item Discovered that question decomposition is not helpful for machine reading but rather harmful. 
%\item Conducted a qualitative error analysis, showing that machine reading using question decomposition struggle due to compound error and question decomposition in bad formats.
%\item Published paper at {\em Proceedings of the Empirical Methods in Natural Language Processing 2023}.
%\end{itemize}


%{\bf Leveraging Multiple Teachers for Test-Time Adaptation of Language-Guided Classifiers} \\
%{\em  The University of North Carolina at Chapel Hill}\hfill{Chapel Hill, NC}\\
%{\em Research Assistant, Mentor: Prof.Shashank Srivastava}\hfill{Apr 2022 - May 2023}
%\begin{itemize}[noitemsep, topsep=0pt,leftmargin=*]
%\item Present a framework for test-time adaptation of language explanation-guided classifiers towards a specific task during inference.
%\item Achieved 8\% higher classification accuracy by utilizing label aggregation with language models, including RoBERTa, T0, OPT, and Flan-T5, for test-time adaptation, and three times better accuracy comparing to baselines with zero-shot learning.
%\item Further improved the accuracy by 20\% by incorporating self-learning by fine-tuning pre-trained LM on noisy labeled data.
%\item Conducted qualitative analysis for framework's interpretability.
%\item Simplify and improve the label aggregation technique by replacing hand-written labeling functions with LM.
%\item Published paper at {\em Findings of the Empirical Methods in Natural Language Processing 2023}.
%\end{itemize}

%{\bf Compositional Generalization for Kinship Prediction through Data Augmentation} \\
%{\em  The University of North Carolina at Chapel Hill}\hfill{Chapel Hill, NC}\\
%{\em Research Assistant, Mentor: Prof.Shashank Srivastava}\hfill{Feb 2021 - Mar 2022}
%\begin{itemize}[noitemsep, topsep=0pt,leftmargin=*]
%\item Evaluated empirically the utility of data augmentation and intermediate structured representations towards compositional generalization for the task of kinship prediction from a story.
%\item Tested the impact of incorporating data augmentation and intermediate structured data on model's performance. Data augmentation boosted generalization performance by around 20\% on average relative to a baseline model from prior work. 
%\item Found that predicting and using intermediate kinship graphs led to a deterioration in the generalization of kinship prediction.
%\item Published paper at {\em Proceedings of the 4th Workshop of Narrative Understanding (WNU2022)}.
%\end{itemize}

%{\bf A Multilingual COVID-19 Question Answering System} \\
%{\em University of California, Santa Barbara}\hfill{Santa Barbara, CA}\\
%{\em Visiting Research Assistant, Mentor: Prof.William Wang}\hfill{May 2021 - Aug 2021}
%\begin{itemize}[leftmargin=*]
%\item Established a multilingual COVID-19 Question Answering system using mBERT and XLM-Roberta models.
%\item Focused primarily on building and training the Reading Comprehension part of the QA system, including collecting %nd processing the data, and performing machine translation on large-scale natural language dataset.
%\item Performed large-scale neural machine translation on entire CORD-19 dataset to acquire data in Chinese and French to alleviate the data scarcity problem in foreign languages.
%\item Reached an F1 score of 60.5 for the reading comprehension model of the final QA system.
%\end{itemize}

%{\bf Temporal and Spatial Analysis on U.S. Domestic Mobility} \\
%{\em Research Assistant}\hfill{Chapel Hill, NC}\\
%{\em Mentor: Prof.Weizi Li at the University of Memphis, Prof.Lei %Lin at University of Rochester}\hfill{Jun 2020 - Oct 2020}
%\begin{itemize}[leftmargin=*]
%\item Investigated the impact of COVID-19 on the United States %mobility through temporal and spatial analysis with network models.
%\item Collected relative datasets from GitHub, Kaggle, New York Times and Google Search Trends; Preprocessed data with Python and implemented the model with Networkx package.
%\item Published paper in the 20th and 21st Joint COTA International Conference of Transportation Professionals.
%\end{itemize}
\end{rSection}


%----------------------------------------------------------------------------------------
%   Publication SECTION
%----------------------------------------------------------------------------------------
\begin{rSection}{Publications}
\begin{itemize}[noitemsep, topsep=0pt, leftmargin=*]
\item {\bf Kangda Wei}, Hao Chen, Pei Chen, Zhe Su, Tianyi Wu, Ankit Patel, Baixiang Huang, Changsheng Wang, Qizhang Feng, Boshi Huang, Jing Huang, Dongmei Jia, Gokce Keskin. SAGE: Self-Attribution-Guided Evolving. %In submission
\item Leyuan Tang, {\bf Kangda Wei}, Tianyu Jiang, Ruihong Huang. Measuring LLM Sycophancy under Sustained Multi-Turn Pressure. %In submission
\item {\bf Kangda Wei}, Ziqian Bi, Ashish Molakalapalli, Haonan Duan, Xiaowei Zou, Vincent Du, Alex Carter, Ruihong Huang. Predicting LLM Training and Inference Runtime: A Survey of Methods, Datasets, and Open Challenges. {\em Findings of the Empirical Methods in Natural Language Processing 2026.}
\item {\bf Kangda Wei}, Ruihong Huang. MMR-GRPO: Accelerating GRPO-Style Training through Diversity-Aware Reward Reweighting. {\em Findings of the Association for Computational Linguistics (ACL) 2026}
\item Cheng Zhang, Rajasekhar Kakarla, {\bf Kangda Wei}, Ruihong Huang. ENG-DRB: PDTB-style Discourse Relation Bank on Engineering Tutorial Video Scripts. {\em Findings of The Asian Federation of Natural Language Processing and The Association for Computational Linguistics 2025}
\item {\bf Kangda Wei}, Hasnat Md Abdullah, Ruihong Huang. Mitigating Gender Bias via Fostering Exploratory Thinking in LLMs. {\em Findings of the Empirical Methods in Natural Language Processing 2025.}
\item Abhilekh Borah, Hasnat Md Abdullah, {\bf Kangda Wei}, Ruihong Huang. CliME: Evaluating Multimodal Climate Discourse on Social Media and the Climate Alignment Quotient (CAQ). {\em NLP4PosImpact Workshop, Association for Computational Linguistics 2025. Best Paper Award.}
\item {\bf Kangda Wei}, Xi Shi, Jonathan Tong, Sai Ramana Reddy, Anandhavelu Natarajan, Rajiv Jain, Aparna Garimella, Ruihong Huang. LegalCore: A Dataset for Event Coreference Resolution in Legal Documents. {\em Findings of Association for Computational Linguistics 2025.}
\item {\bf Kangda Wei}, Zhengyu Zhou, Bingqing Wang, Jun Araki, Lukas Lange, Ruihong Huang, Zhe Feng. Multi-Agent Video Understanding for Advanced Indexing of Lecture Videos. {\em U.S Patent.}
\item Hasnat Md Abdullah, Tian Liu, {\bf Kangda Wei}, Shu Kong, Ruihong Huang. UALBench: The First Comprephensive Unusual Activity Localization Benchmark. {\em Winter Conference on Applications of Computer Vision (WACV) 2025.}
\item {\bf Kangda Wei}, Aayush Guatam, Ruihong Huang. Are LLMs Good Annotators for Discourse-level Event Relation Extraction? {\em Findings of the Empirical Methods in Natural Language Processing 2024.}
\item Jiangshu Du, Yibo Wang, Wenting Zhao, ... {\bf Kangda Wei}, et al. LLMs Assist NLP Researchers: Critique Paper (Meta-)Reviewing. {\em Proceedings of the Empirical Methods in Natural Language Processing 2024.}
\item {\bf Kangda Wei}, Dawn Lawrie, Benjamin Van Durme, Yunmo Chen, Orion Weller. When Do Decompositions Help for Machine Reading? {\em Proceedings of the Empirical Methods in Natural Language Processing 2023}
\item {\bf Kangda Wei}, Sayan Ghosh, Rakesh Menon, and Shashank Srivastava. Leveraging Multiple Teachers for Test-Time Adaptation of Language-Guided Classifiers. {\em Findings of the Empirical Methods in Natural Language Processing 2023}
\item {\bf Kangda Wei}, Sayan Ghosh, and Shashank Srivastava. 2022. Compositional Generalization for Kinship Prediction through Data Augmentation. {\em Proceedings of the 4th Workshop of Narrative Understanding (WNU2022)}, pages 13–19, Seattle, United States. Association for Computational Linguistics.
\item *Songhe Wang, {\bf *Kangda Wei}, Lei Lin, Weizi Li. “Spatial-temporal Analysis of COVID-19's Impact on Human Mobility: the Case of the United States,” in the {\em 20th and 21st Joint COTA International Conference of Transportation Professionals.} *Co-author: equal contribution
\end{itemize}
\end{rSection}

\begin{rSection}{Professional Services}
\begin{itemize}[noitemsep, topsep=0pt, leftmargin=*]
\item Reviewer for ACL Rolling Review (ARR) since February 2024 (covering ACL/NAACL/EMNLP/EACL/AACL).
\item Reviewer for IEEE/CVF Winter Conference on Applications of Computer Vision (WACV), 2025.
\item Reviewer for Neural Information Processing Systems (NeurIPS), 2026.
\end{itemize}
\end{rSection}

\begin{rSection}{Teaching Services}
%{\bf Graduate Teaching Assistant:}
\begin{itemize}[noitemsep, topsep=0pt, leftmargin=*]
\item TA COMP211 System Fundamentals, UNC-Chapel Hill, Department of Computer Science, Spring 2023
\item COMP431 Internet Services and Protocols, UNC-Chapel Hill, Department of Computer Science, Fall 2022
\end{itemize}
\end{rSection}

\begin{rSection}{Awards}
\begin{itemize}[noitemsep, leftmargin=*]
\item Outstanding PhD Research Award (one per year), Texas A\&M University CSE Department, 2026
\item Texas A\&M University Graduate School Travel Grant, 2026
\item Texas A\&M University CSE Department Travel Grant, 2023, 2024, 2025, 2026
\end{itemize}
\end{rSection}


%\end{rSection}
%----------------------------------------------------------------------------------------
%   Projects SECTION
%----------------------------------------------------------------------------------------

%\begin{rSection}{Notable Porjects}
%{\bf Email Plug-in Startup}\hfill{Jan 2022 - May 2022}
%\begin{itemize}[leftmargin=*]
%\item Built a classification model from scratch for email classification using PyTorch and Huggingface Transformer. 
%\item Increased the model accuracy from 0.32 to 0.84 with limited amount of annotated data.
%\end{itemize}

%{\bf Cyber-infrastructure: Web-based Application for Sharing Neural Imaging Data}\hfill{Sep 2020 - Jan 2021}
%\begin{itemize}[leftmargin=*]
%\item Built a website where users can register and login to access the neural imaging data of patients. 
%\item Practiced the basis of front-end and back-end engineering using Python Flask and Firebase. 
%\end{itemize}

%\end{rSection}

%----------------------------------------------------------------------------------------
%	Skills SECTION
%----------------------------------------------------------------------------------------

\begin{rSection}{Skills}
\begin{itemize}[noitemsep, topsep=0pt, leftmargin=*]
\item {\bf Programming Languages / Frameworks:} Python, Java, JavaScript, C, HTML, CSS, MATLAB, R, PyTorch
\item {\bf Software \& Tools:} PyCharm, VS Code, Spyder, Jupyter Notebook, Weights \& Biases,, Docker, Slurm, Git, Tableau, LaTeX, Anaconda, RStudio
\item {\bf LLM Training \& Post-training:} SFT, RLHF, DPO, GRPO, Reward Modeling, Prompt Engineering, Synthetic Data Generation
\item {\bf Cloud \& Infrastructure:} AWS, Microsoft Azure, Amazon Mechanical Turk
\item {\bf Language:} Chinese (native), English (proficient)

\end{itemize}
\end{rSection}

\end{document}
